In this section, apply metagradients to three problems in
machine learning: selecting training data, poisoning training data, and
searching for hyperparameters. In each setting we follow the same recipe:
we frame the task as an optimization problem, modify the learning algorithm
of interest to be \textit{smooth}, then solve by first-order
optimizing with meta-gradients---which we refer to, in a catch-all manner across
algorithms, as metagradient descent (MGD). In particular:
we substantially improve on existing dataset selection methods
(Section \ref{sec:datacomp},
Section \ref{sec:less}), perform the first effective
accuracy-degrading data poisoning attack
(Section \ref{sec:poisoning}), and discover one-cycle learning rate
schedules with MGD
(Section \ref{sec:lrsched}).

\subsection{Selecting multimodal training data}
\label{sec:datacomp}
Curating a training dataset from a mass of unfiltered data is a necessary and
influential step in any large-scale machine learning pipeline. Deciding how to
curate such a dataset is a challenging problem that has attracted substantial
recent
interest~\citep{fang2022data,abbas2023semdedup,engstrom2024dsdm,gadre2024datacomp}.
In this section, we frame pre-training data selection as an optimization
problem, and then solve this problem by first-order optimizing with
metagradients. Applying our method to the DataComp-small benchmark
\citep{gadre2024datacomp}, we greatly improve on the state-of-the-art (our
  improvement over state-of-the-art is roughly the same as the
improvement of state-of-the-art over training on random data).

\subsubsection{Setup}
The goal of {dataset selection} is to choose a training data subset (out of
a broad pool of data) that maximizes trained machine learning model
performance.
Given this goal,
dataset selection has a natural interpretation as a
combinatorial metaparameter optimization problem.
In particular, in the language of Section \ref{subsec:meta_gradient},
for a training set of size $n$, let
\begin{itemize}
  \item[(a)] the metaparameters $\mathbf{c} \in \mathcal{C} :=
    \mathbb{Z}_{\geq 0}^n$ be non-negative data counts representing
    the number of times each training sample repeats in the training data;
  \item[(b)] the algorithm $\mathcal{A}$ be a standard large-scale
    learning procedure, which runs on a training set
    comprising $c_i$ copies of each sample $i$ for $i \in [n]$;
  \item[(c)] the output function $\phi$ be the loss of the trained model on a target
    distribution $\dtarg{}$.
\end{itemize}
Then, defining $f(\mathbf{c}) := \phi(\mathcal{A}(\mathbf{c}))$ (as
in Section~\ref{subsec:meta_gradient}),
our goal is to find the data counts $\mathbf{c^*}$ that solve
\begin{equation}
  \label{eq:ds}
  \mathbf{c}^* := \argmin_{\mathbf{c} \in \mathcal{C}} f(\mathbf{c}).
\end{equation}

\subsubsection{Gradient descent on training data}
\label{sec:gdod}
Metagradients let us \textit{directly} minimize the target task
loss \eqref{eq:ds}
with respect to the choice of training data.
At a high level, our algorithm operates as follows:
we start with a randomly chosen set of training data, then
iteratively update the dataset selection using
metagradients with respect to importance weights
placed on each training datapoint.
The specifics of our method are in Algorithm~\ref{alg:depsing};
we describe its core ideas below.

\paragraph{Idea 1: A surrogate algorithm.} We cannot use metagradients to
optimize \eqref{eq:ds} directly, because the metaparameters of
interest $\mathbf{c}$ are
discrete counts (and so the algorithm $\mathcal{A}$ is non-differentiable with 
respect to $\mathbf{c}$). 
To circumvent this problem, we relax $\mathcal{A}$: we define a surrogate algorithm
$\mathcal{A}_\mathbf{c}'$ that takes in a {\em continuous}
metaparameter $\mathbf{z} \in
\mathbb{R}^n$, whose metagradient we \textit{can} compute, then optimize using
the metagradient on $\mathcal{A}_\mathbf{c}'$.

This surrogate learning algorithm $\mathcal{A}_\mathbf{c}'$ maps a metaparameter 
$\mathbf{z} \in \mathbb{R}^n$ (representing a perturbation to training data weights) 
to a machine learning model.
The surrogate is defined by a set
of counts $\mathbf{c} \in \mathbb{Z}_+^n$, and a hyperparameter $k$ denoting a
specific training iteration, both of which we bake into the surrogate algorithm
itself. 
Given a metaparameter $\mathbf{z} \in \mathbb{R}^n$, the algorithm
$\mathcal{A}_\mathbf{c}'$ trains a model ``as usual'' using the fixed counts
$\mathbf{c}$. That is, it makes $c_i$ copies of each training sample $i$,
shuffles and partitions the data into batches, and then at each iteration
minimizes the batch loss with a step---just as the original learning algorithm
$\mathcal{A}$. At iteration $k$, however, in addition to the original loss on
the $k$-th batch, the algorithm upweights {\em each} training sample $i$
according to the metaparameter $z_i$. In other words, the objective at iteration
$t$ of the surrogate algorithm $\mathcal{A}_\mathbf{c}'$ is
\begin{equation*}
  \label{eq:surrogate_objective}
  \ell'_t(\theta) :=
  \begin{cases}
    \sum_{x \in t^{\text{th}} \text{ batch}} \ell(x; \theta) &
    \!\!\text{if } t \neq k \\
    \sum_{x \in t^{\text{th}} \text{ batch}} \ell(x; \theta) +
    \sum_{i=1}^n z_i \ell(x_i; \theta) &\!\!\text{if } t = k
  \end{cases}
\end{equation*}
where $\ell(x; \theta)$ is the training loss on example $x$.

Observe that when $\mathbf{z} = \mathbf{0}_n$, the algorithm
$\mathcal{A}_\mathbf{c}'$  is identical to the standard learning algorithm
$\mathcal{A}$. And while $\mathcal{A}$ was a function of (nondifferentiable)
discrete data counts $\mathbf{c}$, $\mathcal{A}_\mathbf{c}'$ is differentiable
with respect
its input $\mathbf{z}$, and so we can compute the metagradient
\[
  \mathbf{g} := \nabla_{\mathbf{z}}
  \phi\big(\mathcal{A}_\mathbf{c}'(\mathbf{z})\big)\big\vert_{\mathbf{z}
  = \mathbf{0}_n}.
\]
Intuitively, the entries of the metagradient $\mathbf{g}$ capture the effect of
adding an infinitesimal amount of each training sample $i$ to the training data
at iteration $k$. A positive entry $g_i$ indicates that adding an infinitesimal
amount of sample $i$ to the training data would increase the loss,  and a
negative entry indicates that adding an infinitesimal amount of sample $i$ to
the training data would decrease the loss; the slot at $i$ represents the
(estimated) effect of adding a copy of sample $i$ to the training
data at every batch containing the sample.

\paragraph{Idea 2: Block coordinate descent.}
We then use the metagradient $\mathbf{g}$ to iteratively update our selected
dataset. We update data counts as
\begin{equation}
  \label{eq:update}
  \mathbf{c} \gets \mathbf{c} - \mathrm{sign}(\mathbf{g}) \odot
  \mathbf{m}, \quad \mathbf{m} \sim \mathrm{Bernoulli}(p),
\end{equation}
where $p$ is a hyperparameter controlling the fraction of sample counts to
update. This algorithm resembles a block coordinate descent algorithm
\citep{ortega2000iterative}, with the main difference being that we take signed
gradient steps with step size $1$ (projected onto non-negative integers) to
ensure that the counts remain well-defined. As a result, $p$ implicitly controls
the algorithm's step size.

Applying \eqref{eq:update} concludes a single optimization step.
By repeating this process of estimating the metagradient, updating our
counts vector, then constructing a new training dataset, we iteratively
improve the selected data.  Pseudocode for our algorithm can be found in Algorithm \ref{alg:depsing}.

\begin{algorithm}[h]
    \SetKwFor{For}{\textcolor{forcolor}{for}}{\textcolor{forcolor}{do}}{}
    \SetKw{Return}{\textcolor{forcolor}{Return}}{}
    \SetKwFunction{traversereverse}{\textcolor{fncolor}{reverse\_inorder\_traversal}}
    \DontPrintSemicolon
	\KwIn{initial data counts $\mathbf{c} \in \mathbb{Z}_{\geq 0}^n$, learning algorithm $\mathcal{A}$, output function $\phi$}
    \nonl\hspace{0pt}{\textbf{Hyperparameters:}} step size $p$, \# opt steps $T$, iteration number $k$ \;
    \For{$t \gets 1$ \color{forcolor}{\textbf{to}} \color{black}$T$}{
		$\mathbf{z} \gets \mathbf{0}_n$ \tcp{Build input to surrogate}
		$\mathbf{g} \gets \frac{\partial \phi(\mathcal{A}'_\mathbf{c}(\mathbf{z}))}{\partial \mathbf{z}}$ \tcp{Calculate metagradient using \textsc{Replay}}
		$\mathbf{m} \gets$ sample from $\textrm{Bernoulli}(p)$ \tcp{Sample indices to step on}
		$\mathbf{c} \gets \mathbf{c} - \mathrm{sign}(\mathbf{g}) \odot \mathbf{m}$ \tcp{Take optimization step}
    }
	\Return $\mathbf{c}$ \tcp{Return final data counts}
	\caption{Dataset selection using using metagradient descent (MGD).}
    \label{alg:depsing}
\end{algorithm}

\subsubsection{Results}
We evaluate our data selection algorithm using DataComp
\cite{gadre2024datacomp}, a standardized framework for evaluating data selection
methods for multimodal models. Algorithm~\ref{alg:depsing} greatly
improves on the state-of-the-art for the benchmark. Below, we
describe the setting, outline our
method, and conclude with our results.

\paragraph{Setting.}
DataComp \citep{gadre2024datacomp} is a multimodal model training competition 
and benchmark for evaluating dataset selection methods. 
DataComp provides a \textit{fixed} learning algorithm
chosen in advance by the organizers and a large fixed \emph{candidate pool} of
internet data. The goal is to choose a subset of the
candidate pool---possibly with repeated datapoints---that yields the best-performing
model after training with the given learning algorithm, as measured by a
predetermined set of 38 benchmarks. Given a submission subset, the mean score on
the evaluation datasets for a model trained with that subset is taken as the
final ``score.''  DataComp offers four separate ``scales'' requiring different
amounts of compute; we focus on the \emph{small} scale in this paper
due to compute limitations.

\paragraph{Method.}
We select data with MGD (Algorithm~\ref{alg:depsing}) to minimize loss on
data on a ``target set'' that is distributionally similar to the DataComp
benchmark tasks, and select hyperparameters with a held-out ``validation
set.'' In particular, we construct target and validation sets by taking
samples from the DataComp evaluation tasks with extra samples available beyond
those used in the DataComp test set (e.g., ImageNet, one of the tasks in
DataComp, has a training set in addition to the test set evaluated in DataComp).
See Appendix~\ref{app:clip_details} for the exact details of the target and
validation sets, the precise hyperparameters used with
Algorithm~\ref{alg:depsing}, and a discussion on scalability (including further
engineering details on executing our algorithm efficiently).

\paragraph{Results.}
MGD greatly outperforms the current state-of-the-art: the difference in accuracy
between MGD and the current best method is roughly as large as the difference
between the previous state-of-the-art (EcoDatum~\citep{ecodatum2024ecodatum})
and training on randomly chosen data
(cf.\ Figure~\ref{fig:datacomp_score_over_time}).
Inspecting scores over the course of the optimization in Figure
\ref{fig:datacomp_score_over_time}, we find that only a few steps are necessary
to outperform previous methods.

\begin{figure}[ht]
  \centering
  \begin{subfigure}[t]{0.49\textwidth}
    \raisebox{-\height}{\pgfplotsset{colormap/Set2}
\pgfplotsset{scaled y ticks=false}
\pgfplotsset{grid style={dashed,gray}}

\begin{tikzpicture}
    \begin{axis}[
      xlabel={Metagradient steps},
      ylabel={DataComp Score},
      height={4.8cm},
      width={7.5cm},
      colormap name={Set2},
      ytick distance=0.03,
      xmin=-3,
      xmax=57,
      grid=both
    ]

    \addplot[mark=none, index of colormap=3, domain=-10:60, samples=2, dashed, very thick] {0.1377730386384371};
    \addlegendentry{No filtering}
    \addplot[mark=none, index of colormap=1, domain=-10:60, samples=2, dashed, very thick] {0.173};
    \addlegendentry{Best DataComp baseline}
    \addplot[mark=none,, index of colormap=2, domain=-10:60, samples=2, dashed, very thick] {0.182};
    \addlegendentry{Prev. SOTA (EcoDatum)}
    \addplot [scatter,
                index of colormap=0,
                very thick,
                scatter/@pre marker code/.code={3
                    \def\markopts{index of colormap=0}
                    \expandafter\scope\expandafter[\markopts]
                },
                scatter/@post marker code/.code={
                    \endscope
                },] 
    table[
      x=it,
      y={score},
      col sep=comma,
    ]{pgffigs/datacomp_small_result/datacomp_small_result.csv};
    \addlegendentry{MGD-DS (ours)}

    \legend{}
    \end{axis}
\end{tikzpicture}
}
  \end{subfigure}
  \hfill
  \begin{subfigure}[t]{0.5\textwidth}
    \begin{tabular}[t]{@{}lcc@{}}
\toprule
\textbf{Method}  &  \textbf{Score} & $\Delta$ \\
\midrule
\tikz[baseline=-0.5ex] {\draw [dashed, color=pgfcolor3, very thick] (0,0) -- (0.6,0)} Baseline: No filtering  & 0.13 & -- \\
\tikz[baseline=-0.5ex] {\draw [dashed, color=pgfcolor1, very thick] (0,0) -- (0.6,0)} Best baseline from \cite{gadre2024datacomp}  & 0.17 & \textcolor{mygreen}{+0.04} \\
\tikz[baseline=-0.5ex] {\draw [dashed, color=pgfcolor2, very thick] (0,0) -- (0.6,0)} Previous SOTA \cite{ecodatum2024ecodatum}  & 0.18 & \textcolor{mygreen}{+0.05} \\
\tikz[baseline=-0.5ex] {\draw [color=pgfcolor0, very thick] (0,0) -- (0.6,0); \filldraw [color=pgfcolor0] (0.3,0) circle (0.08);} \textbf{MGD-DS (ours)} & \textbf{0.22} & \textbf{\textcolor{mygreen}{+0.09}} \\
\bottomrule
\end{tabular}

  \end{subfigure}
  \caption{MGD dataset selection greatly outperforms existing methods (improving
      over the previous SOTA by as much as the previous SOTA improves over no
    filtering at all). We compare DataComp scores for MGD (over
    optimization steps),
    training on the entire candidate pool, the best baseline
    originally proposed by
  DataComp, and the previous SOTA~\citep{ecodatum2024ecodatum}.}
  \label{fig:datacomp_score_over_time}
\end{figure}



\subsection{Selecting instruction-tuning data}
\label{sec:less}
In our second application, we select training data for instruction fine-tuning
(IFT) using the same MGD-based method detailed in Algorithm~\ref{alg:depsing} of
Section~\ref{sec:datacomp}. As with multimodal data, training on the ``right''
post-training data (such as the ``right'' IFT data) can greatly impact
deployment-time model
performance~\citep{liu2024deepseek,dubey2024llama,taori2023stanford}. MGD
improves over baselines at choosing IFT data for
MMLU~\citep{hendrycks2021measuring}, a general knowledge task, and
BBH~\citep{suzgun2022challenging}, a reasoning/chain-of-thought task.

To overview this section: we start by detailing the setting, then describe the
specifics of our MGD instantiation before concluding with results.

\begin{figure}[t]
  \centering
  \begin{subfigure}[t]{0.49\textwidth}
    \centering
    \raisebox{-0.605\height}{ \pgfplotsset{colormap/Set2}
\pgfplotsset{scaled y ticks=false}
\pgfplotsset{grid style={dashed,gray}}

\begin{tikzpicture}
\begin{axis}[
    ybar,
    axis lines=box,
    ytick pos=left,
    axis on top=false,
    height=4.8cm,
    width=7.5cm,
    colormap name={Set2},
    ymin=-0.2,
    ymax=1.7,
    ylabel={$\Delta$ Accuracy (\%)},
    symbolic x coords={BBH,MMLU},
    xtick=data,
    xtick style={draw=none},
    enlarge x limits=0.5,
    bar width=15pt,
    grid=major,
    xmajorgrids=false,
    extra y ticks={0},
    extra y tick labels={},
    extra y tick style={
      grid=major,
      major grid style={black, solid}
    },
]

\addplot[
    fill=mapped color,
    draw=mapped color,
    point meta=0,
    index of colormap=1,
] coordinates {
    (BBH,-0.047672)
    (MMLU,0.540605)
};

\addplot[
    fill=mapped color,
    draw=mapped color,
    point meta=0,
    index of colormap=0,
] coordinates {
    (BBH,1.481074)
    (MMLU,1.297280)
};

\end{axis}
\end{tikzpicture}
 }
  \end{subfigure}
  \hfill
  \begin{subfigure}[t]{0.5\textwidth}
    \centering

    
\begin{tabular}{@{}l*{2}{cc}@{}}
	\toprule
	& \multicolumn{2}{c}{\textbf{BBH} \citep{suzgun2022challenging}} & \multicolumn{2}{c}{\textbf{MMLU} \citep{hendrycks2020measuring}} \\
	\cmidrule(lr){2-3} \cmidrule(l){4-5}
	& Acc. & $\Delta$ & Acc. & $\Delta$ \\
	\midrule
	All Data & 35.2\% & $-$ & 41.2\% & $-$ \\
	\textcolor{pgfcolor1}{\rule{0.26cm}{0.26cm}} LESS & 35.2\% & \textcolor{myred}{$-0.0$\%} & 41.8\% & \textcolor{mygreen}{$+0.5$\%} \\
	\textcolor{pgfcolor0}{\rule{0.26cm}{0.26cm}} \textbf{MGD-DS} & \textbf{36.7\%} & \textbf{\textcolor{mygreen}{$\mathbf{+1.5}$\%}} & \textbf{42.5\%} & \textbf{\textcolor{mygreen}{$\mathbf{+1.3}$\%}} \\
	\bottomrule
\end{tabular}

  \end{subfigure}
  \caption{MGD dataset selection outperforms baselines. Comparing to training
    on all the data: it achieves over double the margin of improvement of LESS
    on MMLU, and improves by $+1.5\%$ on BBH (where LESS does not improve at
  all). The $\Delta$ column denotes improvement over not filtering.}
  \label{fig:ift_dss}
\end{figure}

\paragraph{Setting.}
We adopt the setting of LESS~\citep{xia2024less}. Here, the goal is to select a
training data subset from four combined IFT datasets (Flan
  V2~\citep{longpre2023flan}, CoT~\citep{wei2022chain},
  DOLLY~\citep{conover2023free}, and Open Assistant 1
\citep{kopf2024openassistant}) to maximize accuracy on a given target task. We
consider two target tasks from LESS: MMLU (which comprises multiple choice
questions spanning a variety of disciplines) and BBH (a 23 task subset of
BIG-Bench~\citep{srivastava2022beyond}).
In this setup, the data selector can access samples from
each task built from the in-context learning prompts.
Following \citet{xia2024less}, we
fine-tune a $128$-width LoRA~\citep{huang2020dynamics} (in our work, on
Gemma-2B~\citep{team2024gemma}). See Appendix~\ref{app:less} for full details on
the tasks and learning algorithm.

\paragraph{Method.}
We split up the available task samples into two sets---a ``target'' set and a
``validation'' set---then select data with MGD (via Algorithm~\ref{alg:depsing})
by minimizing causal language modeling loss on the ``target'' set of samples.
We select hyperparameters like step size and number of SGD iterations with the
validation set; see Appendix~\ref{app:less} for more details.

\paragraph{Results.}
Comparing with two baselines---training on \textit{all} the data and training
with data selected with LESS~\citep{xia2024less}---MGD yields strictly better
training dataset selections for each target task (cf. Figure~\ref{fig:ift_dss}).
MGD improves most on BBH, a reasoning task, compared to the best
baseline (+1.5\% accuracy). On MMLU, a knowledge-based task, we
outperform baselines by slightly less compared to the best baseline (+0.8\%);
one explanation is that selecting IFT data lends more control over reasoning
than over intrinsic knowledge available in the LM.

Beyond raw accuracy, we inspect losses across each step of the optimization
process. Overall, our method improves validation loss over MGD steps
(cf. Appendix Figures~\ref{fig:bbh_mmlu_over_time}), but also
exhibits signs of overfitting. Given intuition from
overparameterized learning, we might expect this behavior: we optimize a total
of 270,679 ``weights''---each corresponding to a count for a datapoint---to
minimize loss on only a handful of test samples (cf.
Table~\ref{tab:datasets}).

\subsection{Accuracy-degrading (Huber) data poisoning}
\label{sec:poisoning}
The goal of an accuracy-degrading {\em data poisoning} attack is to degrade the
performance of a machine learning model by corrupting a small fraction of its
training data.
Here, the considered threat model is as follows. The attacker is given
a training set $\mathbf{X} = \{x_1, \ldots, x_n\}$
drawn from a distribution $P$,
and a function $\theta(\cdot)$ mapping training data to model
parameters (representing the learning algorithm used by the victim).
The attacker's goal is to return a new training set $\mathbf{X}'$ that differs
from $\mathbf{X}$ in at most $\varepsilon\cdot n$ datapoints while
inducing model
parameters $\theta(\mathbf{X}')$ that perform as poorly as possible on a freshly
drawn test set $T$ from $P$.

Formally, the adversary aims to solve the following optimization problem:
\begin{equation}
  \label{eq:poisoning}
  \arg\max_{\widetilde{x}_1, \ldots, \widetilde{x}_{n_p}}
  \mathbb{E}_{x \sim P}[\ell(x; \theta(\mathbf{X}'))],
\end{equation}
where $\mathbf{X}' = \{\widetilde{x}_1, \ldots, \widetilde{x}_{n_p},
x_{n_p+1}, \ldots, x_n\}$ and $n_p = \lfloor \varepsilon n \rfloor$.
Note that our goal is to degrade the {\em overall} model performance
on a test set $\mathbf{X}_{test}$ drawn from $P$ (in particular,
the test set $\mathbf{X}_{test}$ is {\em unknown} to the adversary).
In this way, this setting resembles
the Huber contamination model in statistics \citep{huber1964robust},
and is strictly more challenging than the usual data poisoning
settings in deep learning (e.g., backdoor attacks
  \citep{gu2017badnets} or attacks that target specific test examples
\citep{koh2017understanding}).

For large-scale machine learning models, finding strong adversaries has proven
challenging---standard loss-minimizing learning algorithms seem quite robust to
maliciously-inserted data~\citep{lu2023exploring}.
In fact, the first non-trivial accuracy degradation data poisoning attacks on
deep models were pioneered by \citet{lu2022indiscriminate} and later improved 
upon by the same set of authors \citep{lu2023exploring}.
Broadly speaking, even constructing attacks that degrade the overall performance of a
learning algorithm by more than the adversarial budget $\varepsilon$
has proven challenging.

\subsubsection{Setup}
We observe that \eqref{eq:poisoning} is a continuous optimization
problem to which we can directly apply our metagradient framework,
approximating the expectation over $P$ by a finite-sample average
over a validation set $\mathbf{X}_{val}$.
In particular, given a (randomly shuffled) training set $\mathbf{X}$
and validation set $\mathbf{X}_{val}$, we set up the
following metaparameter optimization problem (see
Section~\ref{subsec:meta_gradient}):
\begin{itemize}
  \item[(a)] the metaparameter $\mathbf{z} \in \mathcal{X}^{n_p}$ is
    a tensor of $n_p = \lfloor \varepsilon n \rfloor$
    poisoned samples;
  \item[(b)] the algorithm $\mathcal{A}$ maps metaparameters
    $\mathbf{z}$ to a trained model $\mathcal{A}(\mathbf{z})$ by
    replacing the first\footnote{In principle,
      the adversary can also decide {\em which} samples to poison, but
    for simplicity we consider this ``fixed'' case.} $n_p$ samples in
    $\mathbf{X}$ with the samples in $\mathbf{z}$
    and then training on the resulting dataset;
  \item[(c)] the output function $\phi$ evaluates average loss on the
    validation set $\mathbf{X}_{val}$.
\end{itemize}

\subsubsection{Algorithm}
To apply our first-order methods to this problem, we start by
initializing the poisoned data
to be exactly the first $n_p$ samples in $\mathbf{X}$,
\(
  \mathbf{z}^{(0)} := \{\widetilde{x}_i^{(0)} = x_i: i \in [n_p]\}.
\)
Then, for $t = 1, \ldots, T$, we sample a minibatch
$\mathbf{X}_{val}^{(t)}$ from $\mathbf{X}_{val}$
and use \replay{} to compute the metagradient
\[
  \mathbf{g}_t = \frac{d}{d\mathbf{z}} \left(
    \sum_{x \in \mathbf{X}_{val}^{(t)}} \ell(x; \mathcal{A}(\mathbf{z}^{(t-1)}))
  \right),
\]
and update the poisoned data using (projected) gradient ascent:
\[
  \mathbf{z}^{(t)} = \Pi_{\mathcal{X}}\left(
    \mathbf{z}^{(t-1)} + \eta\cdot \text{sign}(\mathbf{g}_t)
  \right),
\]
where $\Pi_{\mathcal{X}}$ is the projection operator onto the sample
space $\mathcal{X}$.
(For example, when $\mathcal{X}$ is the space of image-label pairs,
  $\Pi_{\mathcal{X}}$ clips images' pixel values to $[0, 1]$
and ensures labels are valid probability distributions.)

\begin{figure}[tp]
  \centering
  \includegraphics[width=\textwidth]{figures/cifar10_poisons.png}
  \caption{Examples of poisoned images from Section \ref{sec:poisoning}.}
  \label{fig:poisons}
\end{figure}

\begin{figure}[tp]
  \centering
  \begin{subfigure}[t]{0.48\textwidth}
    \centering
    \raisebox{-\height}{\pgfplotsset{colormap/Set2}
\pgfplotsset{scaled y ticks=false}
\pgfplotsset{grid style={dashed,gray}}

\begin{tikzpicture}
    \begin{axis}[
      xlabel={Metagradient steps},
      ylabel={Test Accuracy},
      height={4.8cm},
      width={7.5cm},
      colormap name={Set2},
      ytick distance=0.03,
      y tick label style={
		/pgf/number format/.cd,
		zerofill,
		precision=2
	  },
      xmin=-15,
      xmax=800,
      grid=both
    ]

    \addplot[mark=none, index of colormap=3, domain=-10:1000, samples=2, dashed, very thick] {0.920};
    \addlegendentry{Original accuracy}
    \addplot[mark=none, index of colormap=1, domain=-10:1000, samples=2, dashed, very thick] {0.912};
    \addlegendentry{Gradient cancelling}
    \addplot [scatter, 
				mark size=0.0,
                index of colormap=0,
                very thick,
                scatter/@pre marker code/.code={3
                    \def\markopts{index of colormap=0}
                    \expandafter\scope\expandafter[\markopts]
                },
                scatter/@post marker code/.code={
                    \endscope
                },] 
    table[
      x=iteration,
      y={test_acc},
      col sep=comma,
    ]{pgffigs/cifar_poisoning/cifar_poisoning.csv};
    \addlegendentry{MGD-DS (ours)}

    \legend{}
    \end{axis}
\end{tikzpicture}
}
  \end{subfigure}
  \hfill
  \begin{subfigure}[t]{0.51\textwidth}
    \centering
    \centering
    \begin{tabular}[t]{@{}lcc@{}}\toprule
      {\bf Model} & {\bf Acc.} & $\Delta$ \\
      \midrule
      \tikz[baseline=-0.5ex] {\draw [dashed, color=pgfcolor3, very
      thick] (0,0) -- (0.6,0)} Original model & 92.0\% & $-$ \\
      \tikz[baseline=-0.5ex] {\draw [dashed, color=pgfcolor1, very
      thick] (0,0) -- (0.6,0)} GradCancel \citep{lu2023exploring} &
      91.2\% & \textcolor{mygreen}{$-$0.80\%} \\
      \tikz[baseline=-0.5ex] {\draw [color=pgfcolor0, very thick]
      (0,0) -- (0.6,0)} \textbf{MGD-DP (ours)} & \textbf{78.1\%} &
      \boldmath{\textbf{\textcolor{mygreen}{$-$13.9\%}}} \\ \midrule
      {\color{darkgray!75} 1-layer NN (for reference)
      \citep{coates2011analysis}} &
      {\color{darkgray!75} 83.3\%} & \textcolor{mygreen}{$-8.7$\%} \\
      \bottomrule
    \end{tabular}
    \label{tab:poisoning}
  \end{subfigure}
  \caption{
    For each iteration of MGD ($x$-axis),
    we train a new model from random
    initialization on a randomly shuffled training set with
    the current iterate of poisoned data injected.
    We evaluate the test accuracy ($y$-axis), and use \replay{} to compute
    the metagradient.
    MGD outperforms the best known attack \citep{lu2023exploring}
    by an order of magnitude and
    (for reference) results in a model that has the same accuracy as a
    single-layer neural network trained on random image features
    \citep{coates2011analysis}.
  }
  \label{fig:poisoning-results}
\end{figure}


\subsubsection{Evaluation}
We use the CIFAR-10 dataset
which consists of 60,000 total images each labeled as one of 10 classes.
We partition the data into 40,000 training examples, 10,000 validation examples,
and 10,000 test examples.
We consider a simple 12-epoch CIFAR-10 training procedure,
which reaches 92.4\% accuracy on the CIFAR-10 test set when applied to the
40,000 training examples. See Appendix
\ref{app:poisoning} for training hyperparameters.

As described above,
we allow the adversary to modify (in-place) a fixed, $\varepsilon$-fraction of
the training data (in our case, 2.5\%)
subject to the constraint that the poisoned images
still lay in the valid (normalized) image range of $[0, 1]$.
We compare our approach---direct optimization of the data poisoning
objective using metagradients---to the state-of-the-art ``Gradient Cancelling''
(GradCancel) method of \citet{lu2023exploring}. In short, GradCancel
is a two-step method which
first finds a poorly performing model,
then finds poisoned data that induces this model as a minimizer of
the training loss.
We present the full method in Appendix \ref{app:poisoning}.

\paragraph{Results.} We find that metagradients enable state-of-the-art
data poisoning attacks, degrading accuracy by 14\%.
In particular, when allowed to corrupt 1000 of the 40,000 training
samples (2.5\%), our method reduces test set accuracy to 78\%---for reference,
the accuracy of a single-layer neural networked trained on the unmodified
CIFAR-10 training set is 83\%.
The strongest existing data poisoning attack,
GradCancel, only reduces test set accuracy by less than 1\%.\footnote{
  \citet{lu2023exploring} report a larger drop; the discrepancy
  is due to our constraint that poisoned data are valid bounded RGB images.
}
In Figure \ref{fig:poisons}, we visualize the
poisoned images and labels found by our method.
In Figure~\ref{fig:poisoning-results},
we visualize the minibatch loss at each step of the
optimization process.

\begin{remark}[Poisoning non-smooth learning algorithms]
  Recall that to apply metagradient descent,
  we alter the learning algorithm $\mathcal{A}$
  to be metasmooth (see Section \ref{sec:smoothness}).
  This involves making modifications such as switching out
  max pooling layers for average pooling layers,
  moving batch normalization layers before activations,
  and scaling down the last layer's output by a factor of 10.
  It is natural to ask: how much does the efficacy of our method
  depend on this smoothness?
  After all, in practice the adversary cannot control the learning algorithm.
  To answer this question, we take the poison samples generated by
  MGD and insert them into the training set of a corresponding
  \smash{\underline{standard}} (i.e., non-metasmooth) learning algorithm.
  We find that our method still
  significantly degrades the performance of the model,
  from $92.8\%$ to $82.6\%$ (a drop of $10.2\%$).
\end{remark}

\subsection{Finding a learning rate schedule}
\label{sec:lrsched}
As a final application, we optimize the learning rate schedule
of stochastic gradient descent (SGD) for training a CIFAR-10 classifier.
By following the metagradients with respect to the learning rate
at each step of training, our procedure matches grid searching
over standard learning rate schedules---despite starting with na\"ive
hyperparameters (a flat learning rate).

Unlike the other applications discussed here,
metagradients do not unlock state-of-the-art performance.
Instead, we discuss this application to illustrate the flexibility
of \replay{}, and in particular its ability to optimize metaparameters
that do not directly affect the loss landscape (i.e., that only affect
the model via the optimization trajectory).
As we discuss in Section \ref{sec:related},
approximate metagradient estimators cannot apply to these metaparameters.

\subsubsection{Setting}
To put learning rate schedule optimization into the metagradient framework,
we parameterize a schedule as a vector $\bm{\eta} \in \mathbb{R}^k$
comprising $k$ evenly-spaced keypoints,
so that the learning rate at iteration $t$ is given by
\begin{equation}
  \label{eq:lr_schedule}
  \eta(t) = \eta_{\lfloor kt/T \rfloor} + \frac{kt/T - \lfloor kt/T
  \rfloor}{\lceil kt/T \rceil - \lfloor kt/T \rfloor} (\eta_{\lceil
  kt/T \rceil} - \eta_{\lfloor kt/T \rfloor}),
\end{equation}
i.e., a linear interpolation between the keypoints.
\begin{itemize}
  \item[(a)] the metaparameter $\bm{\eta} \in \mathbb{R}^k$ is a
    vector of $k$ keypoints;
  \item[(b)] the algorithm $\mathcal{A}$ maps metaparameters
    $\bm{\eta}$ to a trained model $\mathcal{A}(\bm{\eta})$ by
    training a model for $T$ iterations with the learning rate schedule
    defined by \eqref{eq:lr_schedule};
  \item[(c)] the output function $\phi$ evaluates average loss on the
    validation set $\mathbf{X}_{val}$.
\end{itemize}

\subsubsection{Algorithm}
Following the theme of the rest of this section, we optimize the
metaparameter $\bm{\eta}$
directly using MGD. In particular, we initialize the keypoints to be
a flat learning rate schedule,
and then update the keypoints using the metagradient with respect to
the validation loss,
\begin{align*}
  \bm{\eta}^{(t+1)} = \bm{\eta}^{(t)} - \alpha \cdot
  \mathrm{sign}\left(\nabla_{\bm{\eta}}
  \phi(\mathcal{A}(\bm{\eta}^{(t)}))\right).
\end{align*}

\begin{figure}[tp]
    \centering
    \pgfplotsset{colormap/Set2}
\pgfplotsset{scaled y ticks=false}
\pgfplotsset{grid style={dashed,gray}}

\begin{tikzpicture}
    \begin{axis}[
      xlabel={Metagradient steps},
	  ylabel={Accuracy},
      height={5.2cm},
	  width={10cm},
      colormap name={Set2},
	  legend pos=south east,
      legend style={at={(1.1,1)}, anchor=north west, inner sep={6pt}},
      legend cell align={left},
	  xmin=-1,
	  xmax=51,
	  ymin=0.87,
	  ymax=0.95,
      grid=both
    ]
	\addlegendentry{Best grid target acc}
	\addplot[mark=none, index of colormap=3, domain=-2:1000, samples=2, dashed, very thick] {0.9401889};

	\addlegendentry{Best grid test acc}
	\addplot[mark=none, index of colormap=3, domain=-2:1000, samples=2, solid, very thick] {0.9266};

	\addplot[
      index of colormap=0,
	  mark size=0.0,
	  very thick,
	  dashed,
      error bars/.cd,
      y dir=both,
      y explicit,
    ] table[
      x=it,
      y=ymean,
      col sep=comma   %
	]{pgffigs/lr_opt_result/target_over_time.csv};
	\addlegendentry{MGD-LR target acc}

	\addplot[
      index of colormap=0,
	  mark size=0.0,
	  very thick,
      error bars/.cd,
      y dir=both,
      y explicit,
    ] table[
      x=it,
      y=ymean,
      col sep=comma   %
	]{pgffigs/lr_opt_result/test_over_time.csv};
	\addlegendentry{MGD-LR test acc}

	\addplot [
	  name path=upper_target,
	  draw=none,
	  forget plot
	] 
	table [
	  x=it, 
	  y expr=\thisrow{ymean} + \thisrow{yerr}, 
	  col sep=comma
	]{pgffigs/lr_opt_result/target_over_time.csv};

	\addplot [
	  name path=lower_target,
	  draw=none,
	  forget plot
	] 
	table [
	  x=it, 
	  y expr=\thisrow{ymean} - \thisrow{yerr}, 
	  col sep=comma
	]{pgffigs/lr_opt_result/target_over_time.csv};

    \addplot[
      fill opacity=0.2,
	  index of colormap=0,
    ] fill between[
      of=upper_target and lower_target,
    ];

	\addplot [
	  name path=upper_test,
	  draw=none,
	  forget plot
	] 
	table [
	  x=it, 
	  y expr=\thisrow{ymean} + \thisrow{yerr}, 
	  col sep=comma
	]{pgffigs/lr_opt_result/test_over_time.csv};

	\addplot [
	  name path=lower_test,
	  draw=none,
	  forget plot
	] 
	table [
	  x=it, 
	  y expr=\thisrow{ymean} - \thisrow{yerr}, 
	  col sep=comma
	]{pgffigs/lr_opt_result/test_over_time.csv};

    \addplot[
      fill opacity=0.2,
	  index of colormap=0,
    ] fill between[
      of=upper_test and lower_test,
    ];

    \end{axis}
\end{tikzpicture}

    \caption{Target and test accuracies of MGD's learning rate schedule
      over time closely match or exceed those found by a grid search over
      hundreds of combinations of hyperparameters.  95\% confidence
    intervals are plotted for MGD's results.}
    \label{fig:mgd_vs_grid}
\end{figure}

\subsubsection{Evaluation}
We aim to select the learning rate schedule that minimizes the
expected test set loss.
To do so, we reserve 90\% of the CIFAR-10 test set as a ``validation set''
on which we select hyperparameters.
We then use the remaining 10\% as a test set. We compare the
following two approaches:
\begin{itemize}
  \item {\bf Grid search:} We construct a grid over different one cycle
    learning rate schedules, varying the peak learning rate, starting learning
    rate, ending learning rate, and peak learning rate time. In total, we
    consider over 1,000 different learning rate schedules. We use the reserved
    90\% of the test set to select the best learning rate schedule
    from the grid.
  \item {\bf Metagradient descent (MGD):} We run 50 steps of
    MGD starting from a highly suboptimal flat learning rate schedule, aiming to
    minimize loss on the reserved 90\% of the test set. We use the
    last iteration
    of MGD as our learned learning rate schedule.
\end{itemize}
We evaluate the performance of each final learning rate schedule on the held-out
10\% test set and average the results over the same set of 5 unseen
random seeds.
 

\paragraph{Results.}
Comparing our learned hyperparameter schedule to grid search,
as shown in Figure \ref{fig:mgd_vs_grid}, our learned schedule
using only 50 steps of MGD matches the performance of the state-of-the-art
onecycle schedule found via grid search over more than 1000 configurations.
An important caveat, however, is that these numbers are not directly comparable:
grid search can be run in parallel across many machines, while steps
of MGD must be run sequentially.

In practice, we do not advise using MGD for optimizing low-dimensional
hyperparameters, especially ones that have been thoroughly optimized by grid
search (such as CIFAR-10 learning rate schedules
\citep{smith2017super,page2018cifar,li2019exponential,jordan202494}). Still, an
interesting avenue for future work is to study the utility of MGD for optimizing
high-dimensional hyperparameters that are less well-studied, such as
per-parameter/layer learning rates/weight decays for language models, attention
hyperparameters, or gradient preconditioners.

